\chapter{学術英語表現集}
\label{app:english}

論文の各セクションで使える英語表現を整理する。これらの表現は、AIに英文生成を依頼する際の参考にもなるし、AIの出力をチェックする際の基準にもなる。

\section{Introduction で使える表現}

\subsection*{研究背景の記述}

\begin{table}[H]
\centering
\small
\begin{tabular}{p{4.5cm}p{7cm}}
\toprule
日本語の意図 & 英語表現 \\
\midrule
近年、〜が注目されている & In recent years, ... has attracted considerable attention. \\
〜は重要な課題である & ... remains a critical challenge in the field of ... \\
〜の需要が高まっている & There is a growing demand for ... \\
〜は広く研究されてきた & ... has been extensively studied in the literature. \\
〜の進展に伴い & With the rapid advancement of ... \\
〜は〜において重要な役割を果たす & ... plays a crucial role in ... \\
〜はますます重要になっている & ... has become increasingly important. \\
〜の分野では & In the domain of ... \\
〜が実用化されつつある & ... is being deployed in real-world applications. \\
〜に関する研究が活発化している & A growing body of research has focused on ... \\
\bottomrule
\end{tabular}
\end{table}

\subsection*{研究目的の記述}

\begin{table}[H]
\centering
\small
\begin{tabular}{p{4.5cm}p{7cm}}
\toprule
日本語の意図 & 英語表現 \\
\midrule
本研究では〜を提案する & In this paper, we propose ... \\
本研究の目的は〜である & The aim of this study is to ... \\
本研究は〜に取り組む & This work addresses the problem of ... \\
〜を明らかにすることを目指す & We aim to elucidate ... \\
本研究の主な貢献は以下の通りである & The main contributions of this paper are as follows: \\
我々の知る限り、初めて & To the best of our knowledge, this is the first study to ... \\
\bottomrule
\end{tabular}
\end{table}

\subsection*{論文の構成の記述}

\begin{table}[H]
\centering
\small
\begin{tabular}{p{4.5cm}p{7cm}}
\toprule
日本語の意図 & 英語表現 \\
\midrule
本論文は以下のように構成される & The remainder of this paper is organized as follows. \\
第2節では〜について述べる & Section 2 describes ... \\
第3節では〜を紹介する & Section 3 presents ... \\
最後に、第6節で結論を述べる & Finally, Section 6 concludes the paper. \\
\bottomrule
\end{tabular}
\end{table}

\section{Methods で使える表現}

\subsection*{手法の説明}

\begin{table}[H]
\centering
\small
\begin{tabular}{p{4.5cm}p{7cm}}
\toprule
日本語の意図 & 英語表現 \\
\midrule
提案手法の概要を述べる & We provide an overview of the proposed method. \\
〜のアルゴリズムを以下に示す & The algorithm for ... is presented below. \\
本研究では〜を採用した & In this study, we adopted ... \\
〜を〜に適用した & We applied ... to ... \\
〜は以下の3ステップから成る & ... consists of the following three steps. \\
〜に基づいて & Based on ... / Building upon ... \\
〜を拡張して & By extending ... / We extend ... to ... \\
図Xは〜の概要を示す & Figure X illustrates the overview of ... \\
〜と定義する & We define ... as ... \\
〜を仮定する & We assume that ... \\
\bottomrule
\end{tabular}
\end{table}

\subsection*{実験条件の記述}

\begin{table}[H]
\centering
\small
\begin{tabular}{p{4.5cm}p{7cm}}
\toprule
日本語の意図 & 英語表現 \\
\midrule
実験環境は以下の通りである & The experimental setup is as follows. \\
〜を用いて評価を行った & We conducted the evaluation using ... \\
パラメータは〜に設定した & The parameter was set to ... \\
〜を比較手法として用いた & We used ... as a baseline for comparison. \\
実験は〜回繰り返した & The experiment was repeated ... times. \\
〜の条件下で & Under the condition of ... \\
\bottomrule
\end{tabular}
\end{table}

\section{Results で使える表現}

\subsection*{結果の記述}

\begin{table}[H]
\centering
\small
\begin{tabular}{p{4.5cm}p{7cm}}
\toprule
日本語の意図 & 英語表現 \\
\midrule
表Xは〜の結果を示す & Table X shows the results of ... \\
図Xに〜を示す & Figure X presents ... \\
結果から〜が分かる & The results indicate that ... \\
〜が観察された & ... was observed. \\
〜は〜を上回った & ... outperformed ... \\
〜において有意な差が見られた & A significant difference was found in ... \\
〜は〜と一致する & ... is consistent with ... \\
予想通り、〜であった & As expected, ... \\
興味深いことに、〜であった & Interestingly, ... \\
〜は〜\%の改善を示した & ... showed an improvement of ...\% over ... \\
\bottomrule
\end{tabular}
\end{table}

\subsection*{比較の記述}

\begin{table}[H]
\centering
\small
\begin{tabular}{p{4.5cm}p{7cm}}
\toprule
日本語の意図 & 英語表現 \\
\midrule
〜と比較して & Compared with ... / In comparison with ... \\
〜に対して & In contrast to ... / As opposed to ... \\
同様に & Similarly, ... / Likewise, ... \\
一方で & On the other hand, ... / Meanwhile, ... \\
特に & In particular, ... / Notably, ... \\
\bottomrule
\end{tabular}
\end{table}

\section{Discussion/Conclusion で使える表現}

\subsection*{考察}

\begin{table}[H]
\centering
\small
\begin{tabular}{p{4.5cm}p{7cm}}
\toprule
日本語の意図 & 英語表現 \\
\midrule
この結果は〜を示唆している & These results suggest that ... \\
〜の理由として考えられるのは & A possible explanation for ... is that ... \\
この知見は〜と整合する & This finding is in line with ... \\
〜の観点から考えると & From the perspective of ... \\
〜は〜によって説明できる & ... can be attributed to ... \\
さらなる調査が必要である & Further investigation is warranted. \\
\bottomrule
\end{tabular}
\end{table}

\subsection*{限界の記述}

\begin{table}[H]
\centering
\small
\begin{tabular}{p{4.5cm}p{7cm}}
\toprule
日本語の意図 & 英語表現 \\
\midrule
本研究にはいくつかの限界がある & This study has several limitations. \\
本研究の限界として & A limitation of this study is that ... \\
〜は本研究の範囲外である & ... is beyond the scope of this study. \\
データセットの規模が限られている & The dataset size is limited. \\
一般化には注意が必要である & Caution is needed when generalizing ... \\
\bottomrule
\end{tabular}
\end{table}

\subsection*{展望・結論}

\begin{table}[H]
\centering
\small
\begin{tabular}{p{4.5cm}p{7cm}}
\toprule
日本語の意図 & 英語表現 \\
\midrule
今後の研究では〜を検討する & In future work, we plan to investigate ... \\
〜への拡張が期待される & Extension to ... is a promising direction. \\
本研究は〜に貢献する & This study contributes to ... \\
結論として & In conclusion, ... / To summarize, ... \\
本研究で提案した手法は〜に有用である & The proposed method is useful for ... \\
\bottomrule
\end{tabular}
\end{table}

\section{接続表現・論理展開表現}

\subsection*{追加・列挙}

\begin{table}[H]
\centering
\small
\begin{tabular}{ll}
\toprule
機能 & 表現 \\
\midrule
さらに & Furthermore, / Moreover, / In addition, / Additionally, \\
第一に…第二に & First, ... Second, ... / Firstly, ... Secondly, ... \\
同様に & Similarly, / Likewise, / In the same vein, \\
\bottomrule
\end{tabular}
\end{table}

\subsection*{対比・逆接}

\begin{table}[H]
\centering
\small
\begin{tabular}{ll}
\toprule
機能 & 表現 \\
\midrule
しかし & However, / Nevertheless, / Nonetheless, \\
一方で & On the other hand, / Conversely, / In contrast, \\
〜にもかかわらず & Despite ... / In spite of ... / Notwithstanding ...\footnote{やや古風な表現。現代の論文では Despite や In spite of が好まれる。} \\
〜であるが & Although ... / While ... / Whereas ... \\
\bottomrule
\end{tabular}
\end{table}

\subsection*{因果・理由}

\begin{table}[H]
\centering
\small
\begin{tabular}{ll}
\toprule
機能 & 表現 \\
\midrule
したがって & Therefore, / Thus, / Hence, / Consequently, \\
〜のために & Due to ... / Owing to ... / Because of ... \\
その結果 & As a result, / As a consequence, \\
〜を考慮すると & Considering ... / Given that ... \\
\bottomrule
\end{tabular}
\end{table}

\subsection*{例示・具体化}

\begin{table}[H]
\centering
\small
\begin{tabular}{ll}
\toprule
機能 & 表現 \\
\midrule
例えば & For example, / For instance, / e.g., \\
具体的には & Specifically, / More specifically, / In particular, \\
すなわち & That is, / Namely, / i.e., \\
\bottomrule
\end{tabular}
\end{table}

\subsection*{要約・結論}

\begin{table}[H]
\centering
\small
\begin{tabular}{ll}
\toprule
機能 & 表現 \\
\midrule
要するに & In summary, / To summarize, / In brief, \\
結論として & In conclusion, / To conclude, \\
全体として & Overall, / On the whole, / Taken together, \\
\bottomrule
\end{tabular}
\end{table}
